\documentclass[11pt]{article}
\usepackage[localtheorems]{raeez-math-template}

\newtheorem{theorem}{Theorem}
\newtheorem{proposition}[theorem]{Proposition}
\newtheorem{lemma}[theorem]{Lemma}
\newtheorem{problem}[theorem]{Problem}
\newtheorem*{remark}{Remark}

\providecommand{\C}{\mathbb{C}}
\providecommand{\Z}{\mathbb{Z}}
\providecommand{\PP}{\mathbb{P}}
\providecommand{\HH}{\mathbb{H}}
\providecommand{\cA}{\mathcal{A}}
\providecommand{\cH}{\mathcal{H}}
\providecommand{\cI}{\mathcal{I}}
\providecommand{\cO}{\mathcal{O}}
\providecommand{\cS}{\mathcal{S}}
\providecommand{\cT}{\mathcal{T}}
\providecommand{\cV}{\mathcal{V}}
\providecommand{\cW}{\mathcal{W}}
\providecommand{\cY}{\mathcal{Y}}
\providecommand{\fg}{\mathfrak{g}}
\providecommand{\fgl}{\mathfrak{gl}}
\providecommand{\fsu}{\mathfrak{su}}
\providecommand{\CoHA}{\mathrm{CoHA}}
\providecommand{\ChirAlg}{\mathrm{ChirAlg}}
\providecommand{\Borch}{\mathrm{Borch}}
\providecommand{\SpCh}{\mathrm{Sp}^{\mathrm{ch}}}
\providecommand{\PhiFA}{\Phi^{\mathrm{FA}}}
\providecommand{\kcat}{\kappa_{\mathrm{cat}}}
\providecommand{\kch}{\kappa_{\mathrm{ch}}}
\providecommand{\kBKM}{\kappa_{\mathrm{BKM}}}
\providecommand{\kfib}{\kappa_{\mathrm{fiber}}}
\providecommand{\Ntwo}{\mathcal{N}{=}2}

\title{The class-\texorpdfstring{$\cS$}{S} Schur corner of the Igusa denominator algebra}
\author{}
\date{2026-04-30}

\begin{document}
\maketitle

\begin{abstract}
The four-dimensional parent of the \(K3\times E\) Igusa corner is the
class-\(\cS\) theory
\[
  \cT_{24}:=\cT[A_1,\Sigma_{0,24}],
\]
where \(\Sigma_{0,24}\) is a sphere with twenty-four maximal regular
\(\fsu(2)\)-punctures. Its pants decomposition gives
\[
  (n_v,n_h)=(63,88),\qquad c_{4d}=107/6,\qquad c_{2d}=-214.
\]
The closed genus-two curve \(\Sigma_{2,0}\) gives \(c_{4d}=13/6\) and
has no twenty-four puncture labels. The Schur-to-Igusa comparison is the
typed composite
\[
  \chi(\cT_{24})
  \longrightarrow
  \cI_{\mathrm{Schur}}(\cT_{24})
  \xrightarrow{\Pi_{M_{24}}}
  \phi_{0,1}^{K3}
  \xrightarrow{\Borch}
  \Delta_5
  \rightsquigarrow
  \fg_{\Delta_5}.
\]
Here \(D_5=64^{-1}\Delta_5\), \(c_1(0)=10\),
\(\kBKM(\Delta_5)=5\), and \(\Phi_{10}^{\theta}=\Delta_5^2\).
For \(X=K3\times E\), the separated invariants are
\[
  \kcat(X)=0,\qquad \kch(X)=0,\qquad
  \kch^{\mathrm{Heis}}(X)=3,\qquad
  \kBKM(\Delta_5)=5,\qquad \kfib(X)=24.
\]
The native \(\Phi_3\)-output is \(E_1\)-chiral. Corner vertex algebras,
Drinfeld doubles, representation centres, and Schur VOAs may carry
additional two-dimensional structures only after the relevant
construction has been named.
\end{abstract}

\section{The Gaiotto curve}

Let
\[
  \Sigma_{0,24}=(\PP^1;p_1,\ldots,p_{24})
\]
be the \(A_1\) class-\(\cS\) curve with twenty-four maximal regular
\(\fsu(2)\)-punctures, and set
\[
  \cT_{24}:=\cT[A_1,\Sigma_{0,24}].
\]
A pants decomposition of \(\Sigma_{0,n}\) has \(n-2\) trinions \(T_2\)
and \(n-3\) gauged \(SU(2)\) tubes. The \(A_1\) trinion is the free
trifundamental half-hypermultiplet system with effective
\((n_v,n_h)=(0,4)\), and an \(SU(2)\) vector multiplet contributes
\((n_v,n_h)=(3,0)\). Hence
\[
  (n_v,n_h)(\cT[A_1,\Sigma_{0,n}])
  =
  (3(n-3),4(n-2)).
\]
At \(n=24\),
\[
  (n_v,n_h)(\cT_{24})=(63,88).
\]
The Shapere-Tachikawa normalization gives
\[
  c_{4d}=\frac{2n_v+n_h}{12},\qquad
  a_{4d}=\frac{5n_v+n_h}{24},
\]
and the Beem, Lemos, Liendo, Peelaers, Rastelli, and van Rees
Schur-chiral correspondence gives
\[
  c_{2d}=-12c_{4d}.
\]
Therefore
\[
  c_{4d}(\cT_{24})
  =
  \frac{2\cdot63+88}{12}
  =
  \frac{107}{6},
  \qquad
  c_{2d}(\chi(\cT_{24}))=-214.
\]
The Coulomb rank is \(21\), and the flavour algebra is
\(\fsu(2)^{24}\).

\begin{proposition}[The genus-two curve is a different class-\(\cS\) datum]
\label{prop:sigma20-different}
For \(A_1\) on a closed genus-two curve \(\Sigma_{2,0}\), the same
pants arithmetic gives
\[
  (n_v,n_h)=(9,8),\qquad
  c_{4d}=\frac{13}{6},\qquad
  c_{2d}=-26.
\]
The Schur parent with \(c_{4d}=107/6\) is
\(\cT[A_1,\Sigma_{0,24}]\).
\end{proposition}

\begin{proof}
A closed genus-two surface has \(2g-2=2\) trinions and \(3g-3=3\)
gauged \(SU(2)\) tubes. Thus
\[
  (n_v,n_h)=2(0,4)+3(3,0)=(9,8),
\]
so \(c_{4d}=(18+8)/12=13/6\). The punctured sphere has the already
computed vector \((63,88)\), hence central charge \(107/6\). It also
carries the twenty-four puncture labels, the flavour algebra
\(\fsu(2)^{24}\), and the \(M_{24}\)-permutation datum used in the
K3 comparison.
\end{proof}

\section{Schur, Jacobi, and Borcherds data}

The Schur functor assigns to a four-dimensional \(\Ntwo\) SCFT \(\cT\)
a vertex algebra \(\chi(\cT)\). Its graded character is the Schur
index
\[
  \cI_{\mathrm{Schur}}(\cT;q,\mathbf{x})
  =
  \operatorname{Tr}_{\cH_{\mathrm{Schur}}}
  (-1)^Fq^{E-R}\mathbf{x}^{\mathbf f}.
\]
For \(\cT_{24}\), the flavour variables \(\mathbf{x}\) are attached to
the twenty-four maximal \(\fsu(2)\)-punctures. Marking these punctures
by the singular fibres of an elliptic \(K3\) gives an \(M_{24}\)-equivariant
comparison map
\[
  \Pi_{M_{24}}\colon
  \cI_{\mathrm{Schur}}(\cT_{24})
  \longrightarrow
  \phi_{0,1}^{K3}(\tau,z),
\]
where the Gritsenko-Nikulin normalization is
\[
  \phi_{0,1}^{K3}(\tau,z)=y+10+y^{-1}+O(q).
\]
The geometric K3 elliptic genus is \(2\phi_{0,1}^{K3}\). The primitive
Borcherds denominator for \(\fg_{\Delta_5}\) uses
\(\phi_{0,1}^{K3}\), so
\[
  c_1(0)=[q^0y^0]\phi_{0,1}^{K3}=10,\qquad
  \kBKM(\Delta_5)=\frac{c_1(0)}{2}=5.
\]

\begin{proposition}[Igusa normalizations]
\label{prop:igusa-normalizations}
Let \(Z=\begin{psmallmatrix}\tau&z\\ z&\sigma\end{psmallmatrix}\in\HH_2\).
In the theta normalization,
\[
  [q^{1/2}r^{1/2}s^{1/2}]\Delta_5=64.
\]
The monic Borcherds product is
\[
  D_5=64^{-1}\Delta_5,
  \qquad
  \operatorname{den}(\fg_{\Delta_5})=D_5(2Z)=64^{-1}\Delta_5(2Z).
\]
The square normalization gives the usual Igusa line
\[
  \Phi_{10}^{\theta}=\Delta_5^2.
\]
In the Oberdieck-Pixton monic scalar convention,
\[
  \Phi_{10}^{\mathrm{OP}}=D_5^2=4096^{-1}\Delta_5^2.
\]
\end{proposition}

\begin{proof}
The first coefficient fixes the passage from the theta-normalized
Igusa square root to the monic product. Squaring \(D_5\) gives
\(4096^{-1}\Delta_5^2\), while squaring the theta-normalized
\(\Delta_5\) gives the classical weight-ten Igusa form. The weight of
the primitive denominator remains
\(\kBKM(\Delta_5)=c_1(0)/2=5\); the square has weight \(10\).
\end{proof}

\begin{remark}
The symbols in the displayed composite have different types. The Schur
index is a protected character of a four-dimensional theory. The Jacobi
form \(\phi_{0,1}^{K3}\) is a two-variable modular object. The form
\(\Delta_5\) is a genus-two Siegel form. The algebra
\(\fg_{\Delta_5}\) is the generalized Borcherds-Kac-Moody
superalgebra whose signed root supermultiplicities are supplied by the
Fourier coefficients of \(\phi_{0,1}^{K3}\).
\end{remark}

\section{\texorpdfstring{$K3\times E$}{K3xE} invariants}

Let \(X=K3\times E\). The compact Hodge column is
\[
  h^{0,\bullet}(X)=(1,1,1,1),
\]
so the compact Hodge supertrace gives
\[
  \kch(X)=\sum_{q=0}^{3}(-1)^qh^{0,q}(X)=1-1+1-1=0.
\]
The categorical Euler invariant is Kunneth-multiplicative:
\[
  \kcat(X)=\chi(\cO_{K3})\chi(\cO_E)=2\cdot0=0.
\]
The Heisenberg-Mukai specialization has
\[
  \kch^{\mathrm{Heis}}(X)=3,
\]
and the K3 fibre contributes the Mukai lattice rank
\[
  \kfib(X)=\operatorname{rk}\widetilde H(K3,\Z)=1+22+1=24.
\]
Together with Proposition~\ref{prop:igusa-normalizations},
\[
  \{\kcat(X),\kch(X),\kch^{\mathrm{Heis}}(X),
  \kBKM(\Delta_5),\kfib(X)\}
  =
  \{0,0,3,5,24\},
\]
with the two zeros distinguished by construction.

\begin{theorem}[Separated numerical layers]
\label{thm:separated-numerical-layers}
For \(X=K3\times E\), the total-space categorical invariant, compact
chiral Hodge supertrace, Heisenberg-Mukai specialization, Borcherds
weight, and fibre rank are distinct constructions:
\[
  \kcat(X)=0,\qquad
  \kch(X)=0,\qquad
  \kch^{\mathrm{Heis}}(X)=3,\qquad
  \kBKM(\Delta_5)=5,\qquad
  \kfib(X)=24.
\]
No additive formula among these five quantities is part of the
construction.
\end{theorem}

\begin{proof}
The first two equalities are direct Hodge and Kunneth calculations on
the compact total space. The value \(3\) belongs to the
Heisenberg-Mukai specialization \(\SpCh_{K3,E}\circ\PhiFA_3\). The
value \(5\) is the Borcherds weight \(c_1(0)/2\) of the primitive
denominator. The value \(24\) is the rank of the Mukai lattice of the
K3 fibre. These constructions use different objects and different
functorial inputs, so an additive relation would have to supply an
extra map between them. No such map occurs in the Schur-Jacobi-Borcherds
chain or in the compact CY-to-chiral assignment.
\end{proof}

\section{The native \texorpdfstring{$\Phi_3$}{Phi3} output}

For a Calabi-Yau threefold, the native output of the CY-to-chiral
assignment is \(E_1\)-chiral. In the present case the constructed
object has the form
\[
  A_X
  =
  \SpCh_{K3,E}\bigl(\PhiFA_3(D^b\operatorname{Coh}(K3\times E))\bigr)
  \in E_1\text{-}\ChirAlg(E).
\]
An \(E_2\)-structure may appear on a named centre, double, module
category, or external VOA corner layer:
\[
  Z^{\mathrm{der}}_{\mathrm{ch}}(A_X),\qquad
  D(\CoHA),\qquad
  Z(\operatorname{Rep}(A_X)),\qquad
  \chi(\cT_{24}).
\]
Each expression is a different construction. The symbol \(A_X\) denotes
the \(E_1\)-chiral output.

\section{CoHA, Yangian, and \texorpdfstring{$\cW_{1+\infty}$}{W1+infty}}

The toric local model fixes the positive-half convention:
\[
  \CoHA(\C^3)\cong Y^+(\widehat{\fgl}_1).
\]
The full affine Yangian and the \(\cW_{1+\infty}\) vacuum module enter
after passing to the Drinfeld double, completing, and evaluating the
centre:
\[
  D\bigl(Y^+(\widehat{\fgl}_1)\bigr)
  =
  Y(\widehat{\fgl}_1)
  \curvearrowright
  \cW_{1+\infty}\text{-vac}.
\]
For \(K3\times E\), a compact critical CoHA with orientation data and a
Hall-Borcherds recognition map is an additional input. The scalar
identity \(D_5=64^{-1}\Delta_5\) does not supply that compact source
algebra.

\section{Corner vertex algebra boundary}

The Gaiotto-Rapcak algebras \(\cY_{L,M,N}(\psi)\) are vertex algebras
attached to corners of \(GL\)-twisted four-dimensional
\(\mathcal{N}{=}4\) gauge theory. Their finite triples record boundary
rank data. The Igusa denominator requires the rank-three lattice
\(\Lambda^{2,1}_{II}\), the \(\phi_{0,1}^{K3}\)-graded imaginary
root cone, and a genus-two Siegel variable. Thus a finite
\(\cY_{L,M,N}\) has the wrong grading for the full denominator
of \(\fg_{\Delta_5}\).

The Mukai block \(\cY_{4,20,0}\) is a useful finite boundary model:
the numbers \(4\) and \(20\) remember the even and odd Mukai signature
split. It supplies a corner truncation, while the Borcherds imaginary
tail lives in the Siegel denominator layer. Miki's automorphism gives
the local triality of toroidal charts and exchanges horizontal and
vertical affine subalgebras in the completed \(\cW_{1+\infty}\)
environment. Its compact \(K3\times E\) use requires compatibility with
the Hall-Borcherds centre and the \(\Delta_5\) orientation character.

\begin{problem}[Corner avatar with correct inputs]
\label{prob:corner-avatar-correct-inputs}
Construct a sheaf of completed corner vertex algebras
\(\cV_{\mathrm{corner}}\) over the paramodular Humbert locus with the
following data:
\begin{enumerate}[leftmargin=2em]
  \item a finite Mukai-block restriction whose local model is
  \(\cY_{4,20,0}(\psi)\);
  \item a \(\Lambda^{2,1}_{II}\)-graded completion carrying the
  \(\phi_{0,1}^{K3}\)-Fourier coefficients as signed imaginary-root
  supermultiplicities;
  \item a Miki triality compatible with the Weyl \(S_3\)-symmetry of
  the three real simple roots of \(\fg_{\Delta_5}\);
  \item a BRST or centre reduction whose denominator is
  \(D_5(2Z)=64^{-1}\Delta_5(2Z)\);
  \item a comparison map from the Schur VOA \(\chi(\cT_{24})\) through
  the \(M_{24}\)-averaged Jacobi form to this denominator algebra.
\end{enumerate}
Solving this problem would identify a Gaiotto-corner avatar of
\(\fg_{\Delta_5}\). The numerical input from the four-dimensional parent
would remain
\[
  (n_v,n_h,c_{4d},c_{2d})=(63,88,107/6,-214),
\]
and the Borcherds weight would remain \(\kBKM(\Delta_5)=5\).
\end{problem}

\section{Typed register}

\begin{center}
\renewcommand{\arraystretch}{1.25}
\begin{tabular}{lll}
\hline
Layer & Object & Numerical datum \\
\hline
Gaiotto curve & \(\Sigma_{0,24}\) & \(24\) maximal punctures \\
Class-\(\cS\) theory & \(\cT[A_1,\Sigma_{0,24}]\) & \((n_v,n_h)=(63,88)\) \\
Schur VOA & \(\chi(\cT_{24})\) & \(c_{2d}=-214\) \\
Jacobi seed & \(\phi_{0,1}^{K3}\) & \(c_1(0)=10\) \\
Borcherds denominator & \(D_5=64^{-1}\Delta_5\) & \(\kBKM=5\) \\
Igusa square & \(\Phi_{10}^{\theta}=\Delta_5^2\) & weight \(10\) \\
CY threefold & \(K3\times E\) & \(\kcat=\kch=0\) \\
Heisenberg specialization & \(\SpCh_{K3,E}\circ\PhiFA_3\) & \(\kch^{\mathrm{Heis}}=3\) \\
K3 fibre & \(\widetilde H(K3,\Z)\) & \(\kfib=24\) \\
Toric positive half & \(\CoHA(\C^3)\) & \(Y^+(\widehat{\fgl}_1)\) \\
\hline
\end{tabular}
\end{center}

\begin{thebibliography}{99}
\bibitem{Gaiotto2012}
D. Gaiotto, \emph{\(N=2\) dualities}, JHEP 1208 (2012) 034,
arXiv:0904.2715.

\bibitem{ChacaltanaDistler2010}
O. Chacaltana and J. Distler, \emph{Tinkertoys for Gaiotto duality},
JHEP 1011 (2010) 099, arXiv:1008.5203.

\bibitem{ShapereTachikawa2008}
A. Shapere and Y. Tachikawa, \emph{Central charges of
\(\Ntwo\) superconformal field theories in four dimensions},
JHEP 0809 (2008) 109, arXiv:0804.1957.

\bibitem{BLLPRvR2015}
C. Beem, M. Lemos, P. Liendo, W. Peelaers, L. Rastelli, and
B. C. van Rees, \emph{Infinite chiral symmetry in four dimensions},
Commun. Math. Phys. 336 (2015) 1359-1433, arXiv:1312.5344.

\bibitem{GritsenkoNikulin1997}
V. A. Gritsenko and V. V. Nikulin, \emph{Siegel automorphic form
corrections of some Lorentzian Kac-Moody Lie algebras},
Amer. J. Math. 119 (1997) 181-224, arXiv:alg-geom/9504006.

\bibitem{Gritsenko1999}
V. A. Gritsenko, \emph{Elliptic genus of Calabi-Yau manifolds and
Jacobi and Siegel modular forms}, Algebra i Analiz 11 (1999) 100-125,
arXiv:math/9906190.

\bibitem{Borcherds1998}
R. E. Borcherds, \emph{Automorphic forms with singularities on
Grassmannians}, Invent. Math. 132 (1998) 491-562,
arXiv:alg-geom/9609022.

\bibitem{GaiottoRapcak2018}
D. Gaiotto and M. Rapcak, \emph{Vertex algebras at the corner},
JHEP 1901 (2019) 160, arXiv:1703.00982.

\bibitem{ProchazkaRapcak2018}
T. Prochazka and M. Rapcak, \emph{Webs of \(W\)-algebras},
JHEP 1811 (2018) 109, arXiv:1711.06888.

\bibitem{SchiffmannVasserot2013}
O. Schiffmann and E. Vasserot, \emph{Cherednik algebras, W-algebras
and the equivariant cohomology of the moduli space of instantons on
\(A^2\)}, Publ. Math. IHES 118 (2013) 213-342.
\end{thebibliography}

\end{document}
